\documentclass[11pt,a4paper]{article}

\usepackage[utf8]{inputenc}
\usepackage[T1]{fontenc}
\usepackage[spanish,es-noshorthands]{babel}
\usepackage[margin=1.5cm,top=1.5cm,bottom=1.5cm]{geometry}
\usepackage{booktabs}
\usepackage{tabularx}
\usepackage{xcolor}
\usepackage{listings}
\usepackage{enumitem}
\usepackage{titlesec}
\usepackage{hyperref}
\usepackage{graphicx}
\usepackage{tikz}
\usepackage{tcolorbox}
\usepackage{multicol}

\usetikzlibrary{arrows.meta,positioning,shapes.geometric,fit,backgrounds,calc}

\definecolor{primary}{HTML}{1B2A4A}
\definecolor{secondary}{HTML}{2E86AB}
\definecolor{accent}{HTML}{E8651A}
\definecolor{lightbg}{HTML}{F5F7FA}
\definecolor{codebg}{HTML}{F0F2F5}
\definecolor{darktext}{HTML}{2C3E50}
\definecolor{midgray}{HTML}{7F8C8D}
\definecolor{codegreen}{HTML}{27AE60}
\definecolor{codeblue}{HTML}{2980B9}
\definecolor{codegray}{HTML}{95A5A6}
\definecolor{uartcolor}{HTML}{E74C3C}
\definecolor{i2ccolor}{HTML}{27AE60}
\definecolor{usbcolor}{HTML}{2980B9}

\tcbuselibrary{skins,breakable}

\newtcolorbox{infobox}[1][]{
  colback=lightbg,
  colframe=secondary!60,
  fonttitle=\bfseries\small,
  boxrule=0.6pt,
  arc=3pt,
  left=6pt,right=6pt,top=4pt,bottom=4pt,
  #1
}

\lstdefinestyle{pystyle}{
  language=Python,
  backgroundcolor=\color{codebg},
  basicstyle=\ttfamily\footnotesize,
  keywordstyle=\color{codeblue}\bfseries,
  commentstyle=\color{codegreen}\itshape,
  stringstyle=\color{accent},
  numbers=left,
  numberstyle=\tiny\color{codegray},
  numbersep=5pt,
  breaklines=true,
  frame=l,
  framerule=2pt,
  rulecolor=\color{codeblue!50},
  showstringspaces=false,
  tabsize=4,
  xleftmargin=2em,
  aboveskip=0.5em,
  belowskip=0.5em,
}

\lstdefinelanguage{Rust}{
  morekeywords={let,mut,fn,use,extern,crate,mod,pub,impl,Self,match,Ok,Err,
    unwrap,expect,return,if,else,for,in,while,loop,break,continue,as,ref,
    struct,enum,type,const,static,trait,where,async,await,move,true,false},
  sensitive=true,
  morecomment=[l]{//},
  morecomment=[s]{/*}{*/},
  morestring=[b]",
}

\lstdefinestyle{ruststyle}{
  language=Rust,
  backgroundcolor=\color{codebg},
  basicstyle=\ttfamily\footnotesize,
  keywordstyle=\color{codeblue}\bfseries,
  commentstyle=\color{codegreen}\itshape,
  stringstyle=\color{accent},
  numbers=left,
  numberstyle=\tiny\color{codegray},
  numbersep=5pt,
  breaklines=true,
  frame=l,
  framerule=2pt,
  rulecolor=\color{uartcolor!50},
  showstringspaces=false,
  tabsize=4,
  xleftmargin=2em,
  aboveskip=0.5em,
  belowskip=0.5em,
  morekeywords={let,mut,fn,use,extern,crate,mod,pub,impl,Self,match,Ok,Err,unwrap,expect},
}

\lstdefinestyle{tomlstyle}{
  backgroundcolor=\color{codebg},
  basicstyle=\ttfamily\footnotesize,
  breaklines=true,
  frame=l,
  framerule=2pt,
  rulecolor=\color{midgray},
  showstringspaces=false,
  xleftmargin=2em,
  aboveskip=0.5em,
  belowskip=0.5em,
}

\titleformat{\section}{\Large\bfseries\color{primary}}{}{0em}{}[{\titlerule[1pt]\color{secondary}}]
\titleformat{\subsection}{\large\bfseries\color{secondary}}{}{0em}{}
\titleformat{\subsubsection}{\normalsize\bfseries\color{darktext}}{}{0em}{}

\setlength{\parskip}{0.3em}
\setlength{\parindent}{0pt}

\hypersetup{
  colorlinks=true,
  linkcolor=secondary,
  urlcolor=accent,
}

\begin{document}

% ============================================================
% HEADER
% ============================================================

\begin{center}
\begin{tikzpicture}
  \fill[primary] (0,0) rectangle (\textwidth,2.2);
  \node[anchor=west, text=white, font=\LARGE\bfseries] at (0.5,1.3) {Mapa de Buses de Comunicaci\'{o}n};
  \node[anchor=west, text=secondary!40!white, font=\large] at (0.5,0.6) {TurboPi Hiwonder --- Robot de F\'{u}tbol Aut\'{o}nomo};
  \node[anchor=east, text=white!60, font=\small] at (\textwidth-0.5,0.3) {Fuente: \texttt{hardware.py} $\cdot$ Abril 2026};
\end{tikzpicture}
\end{center}

\vspace{0.5em}

% ============================================================
% DIAGRAMA DE TOPOLOGIA
% ============================================================

\begin{center}
\resizebox{0.96\textwidth}{!}{%
\begin{tikzpicture}[
  x=1cm,
  y=1cm,
  rpi/.style={rectangle, draw=primary, fill=primary!12, very thick, rounded corners=5pt, minimum width=3.0cm, minimum height=4.6cm, font=\normalsize\bfseries, align=center},
  busnode/.style={rectangle, draw=#1, fill=#1!10, thick, rounded corners=4pt, minimum width=2.9cm, minimum height=1.3cm, font=\small, align=center},
  groupnode/.style={rectangle, draw=#1!75, fill=#1!6, rounded corners=3pt, minimum width=4.2cm, minimum height=1.6cm, font=\scriptsize, align=left, text width=3.9cm, inner sep=6pt},
  buslabel/.style={font=\scriptsize\bfseries, text=#1, fill=white, inner sep=1pt},
  arr/.style={-{Stealth[length=6pt,width=4pt]}, very thick, color=#1},
]

\node[rpi] (rpi) at (0,0) {\textcolor{primary}{Raspberry Pi}\\[2pt]{\small TurboPi}};

\node[busnode=uartcolor] (uart) at (5.2,2.1)
  {\textcolor{uartcolor}{\textbf{UART}}\\[1pt]\scriptsize\texttt{/dev/ttyAMA0}\\[-2pt]\scriptsize 1,000,000 baud};
\node[busnode=i2ccolor] (i2c) at (5.2,0)
  {\textcolor{i2ccolor}{\textbf{I\textsuperscript{2}C}}\\[1pt]\scriptsize Bus 1 $\cdot$ \texttt{/dev/i2c-1}};
\node[busnode=usbcolor] (usb) at (5.2,-2.1)
  {\textcolor{usbcolor}{\textbf{USB}}\\[1pt]\scriptsize\texttt{/dev/video*}};

\coordinate (rpi_uart_port) at ($(rpi.east)+(0,1.5)$);
\coordinate (rpi_i2c_port) at ($(rpi.east)+(0,0)$);
\coordinate (rpi_usb_port) at ($(rpi.east)+(0,-1.5)$);

\draw[arr=uartcolor] (rpi_uart_port) -- node[buslabel=uartcolor, above]{tx/rx} (uart.west);
\draw[arr=i2ccolor] (rpi_i2c_port) -- node[buslabel=i2ccolor, above]{SDA/SCL} (i2c.west);
\draw[arr=usbcolor] (rpi_usb_port) -- node[buslabel=usbcolor, above]{USB 2.0} (usb.west);

\node[groupnode=uartcolor] (uartgroup) at (9.7,2.1)
  {\textbf{Actuadores UART}\\
   Servo Pan (ID=2, 0--180\textdegree)\\
   Servo Tilt (ID=1, 0--180\textdegree)\\
   Motores M0..M3 (float32)};

\node[groupnode=i2ccolor] (i2cgroup) at (9.7,0)
  {\textbf{Sensores I\textsuperscript{2}C}\\
   Ultras\'{o}nico (0x77, 0--5000\,mm, LED)\\
   Line Follower (0x78, 4x IR)};

\node[groupnode=usbcolor] (usbgroup) at (9.7,-2.1)
  {\textbf{Visi\'{o}n USB}\\
   C\'{a}mara (\texttt{/dev/video*})\\
   Detecci\'{o}n HSV + YOLO 26 nano ONNX};

\draw[arr=uartcolor] (uart.east) -- (uartgroup.west);
\draw[arr=i2ccolor] (i2c.east) -- (i2cgroup.west);
\draw[arr=usbcolor] (usb.east) -- (usbgroup.west);

\end{tikzpicture}%
}
\end{center}

\vspace{0.3em}

% ============================================================
% TABLA RESUMEN
% ============================================================

\begin{infobox}[title={Resumen de Buses}]
\centering\small
\begin{tabularx}{\textwidth}{l l l l X}
\toprule
\textbf{Bus} & \textbf{Dispositivo} & \textbf{Direcci\'{o}n} & \textbf{Baud/Reg} & \textbf{Funci\'{o}n} \\
\midrule
\textcolor{uartcolor}{\textbf{UART}} & SerialBus & \texttt{/dev/ttyAMA0} & 1,000,000 & 2 servos pan/tilt + 4 motores DC omnidireccionales \\
\textcolor{i2ccolor}{\textbf{I\textsuperscript{2}C}} & Ultras\'{o}nico & \texttt{0x77} & Reg \texttt{0x02} LED & Sensor distancia 0--5000\,mm + LED RGB \\
\textcolor{i2ccolor}{\textbf{I\textsuperscript{2}C}} & Line Follower & \texttt{0x78} & Reg \texttt{0x01} & 4 sensores IR digitales (bitmask) \\
\textcolor{usbcolor}{\textbf{USB}} & C\'{a}mara & \texttt{/dev/video*} & --- & Detecci\'{o}n HSV adaptativa + YOLO 26 nano ONNX \\
\bottomrule
\end{tabularx}
\end{infobox}

\newpage

% ============================================================
% DETALLE UART
% ============================================================

\section{\textcolor{uartcolor}{UART} --- \texttt{/dev/ttyAMA0}}

\subsection{Protocolo binario}

Todas las tramas usan header fijo \texttt{0xAA 0x55}, seguido de tipo, longitud, datos y CRC8 (tabla de 256 bytes).

\begin{center}\small
\begin{tabular}{llp{8cm}}
\toprule
\textbf{Campo} & \textbf{Bytes} & \textbf{Descripci\'{o}n} \\
\midrule
Header & \texttt{0xAA 0x55} & Sync bytes fijos \\
Tipo & 1 & \texttt{0x04}=servo, \texttt{0x03}=motor \\
Longitud & 1 & Cantidad de bytes de datos \\
Datos & N & Ver estructura abajo \\
CRC8 & 1 & Sobre campos tipo+longitud+datos \\
\bottomrule
\end{tabular}
\end{center}

\begin{infobox}[title={Estructura de trama Servo (tipo \texttt{0x04})}]
\small
\texttt{dur\_ms(2B LE)} $+$ \texttt{count(1B)} $+$ \texttt{[id(1B) $+$ pos\_pulse(2B LE)]} $\times$ count. \\[3pt]
Servo Pan (ID=2): centro=86\textdegree, rango 0--180\textdegree, \textbf{invertido}. Pulso: $p = 500 + (\theta/180)\times 2000\;\mu s$. \\
Servo Tilt (ID=1): centro=129\textdegree, rango 0--180\textdegree, \textbf{invertido}. Mismo c\'{a}lculo de pulso.
\end{infobox}

\begin{infobox}[title={Estructura de trama Motor (tipo \texttt{0x03})}]
\small
\texttt{0x05(1B)} $+$ \texttt{count(1B=4)} $+$ \texttt{[motor\_id(1B) $+$ speed\_float32(4B LE)]} $\times$ 4. \\[3pt]
Motor 0 = delantero-izq, Motor 1 = delantero-der, Motor 2 = trasero-izq, Motor 3 = trasero-der. \\
Velocidad: \texttt{float32} con signo (negativo=atr\'{a}s). Mapeo diferencial: $(-v_L,\;v_R,\;-v_L,\;v_R)$ con cap=250.
\end{infobox}

% ============================================================
% DETALLE I2C
% ============================================================

\section{\textcolor{i2ccolor}{I\textsuperscript{2}C} --- Bus 1}

Ambos sensores comparten el bus con un \texttt{threading.Lock} para acceso seguro.

\subsection{Sensor Ultras\'{o}nico (\texttt{0x77})}

\begin{itemize}[nosep,leftmargin=1.5em]
  \item \textbf{Distancia:} write \texttt{[0x00]}, luego read 2 bytes como \texttt{uint16\_little\_endian}. Rango 0--5000\,mm. Error $\rightarrow$ 9999.
  \item \textbf{LED RGB:} registro \texttt{0x02}. Datos: \texttt{[mode, R, G, B, R, G, B]}. \texttt{mode=0x00} s\'{o}lido, \texttt{0x01} parpadeo.
\end{itemize}

\subsection{Line Follower (\texttt{0x78})}

\begin{itemize}[nosep,leftmargin=1.5em]
  \item \textbf{Lectura:} registro \texttt{0x01}, 1 byte. Bitmask: bit0=IR1, bit1=IR2, bit2=IR3, bit3=IR4.
  \item \textbf{Calibraci\'{o}n:} 20 muestras para establecer baseline (pasto dominante = \texttt{True}).
  \item \textbf{Detecci\'{o}n:} $\geq$2 sensores cambian respecto al baseline $\Rightarrow$ l\'{i}nea blanca detectada.
\end{itemize}

% ============================================================
% DETALLE CAMARA
% ============================================================

\section{\textcolor{usbcolor}{C\'{a}mara USB}}

\begin{itemize}[nosep,leftmargin=1.5em]
  \item Resoluci\'{o}n 320$\times$240, buffer=1, auto-exposici\'{o}n deshabilitada, exposici\'{o}n manual 150--450 (default 200).
  \item \textbf{HSV adaptativo:} H naranja = 0--10 y 168--179, S$\geq$230, V$\geq$20. Gamma din\'{a}mico 1.0--2.0. Modo estricto/relajado/re-adquisici\'{o}n.
  \item \textbf{Filtros:} area$>$700px, radio$>$15px, circularidad$>$0.20. Centro de hue adaptativo con EMA $\alpha$=0.15.
  \item \textbf{YOLO 26 nano ONNX:} modelo ONNX 320$\times$320, confianza$\geq$0.35, clase 0 = pelota.
\end{itemize}

\newpage

% ============================================================
% EJEMPLOS PYTHON
% ============================================================

\section{Ejemplos Python}

\subsection{UART: enviar comando servo + motor}

\begin{lstlisting}[style=pystyle]
import struct, serial

CRC8 = [0,94,188,226,97,63,221,131,194,156,126,32,163,
  253,31,65,157,195,33,127,252,162,64,30,95,1,227,189,
  62,96,130,220,35,125,159,193,66,28,254,160,225,191,
  93,3,128,222,60,98,190,224,2,92,223,129,99,61,124,
  34,192,158,29,67,161,255,70,24,250,164,39,121,155,
  197,132,218,56,102,229,187,89,7,219,133,103,57,186,
  228,6,88,25,71,165,251,120,38,196,154,101,59,217,
  135,4,90,184,230,167,249,27,69,198,152,122,36,248,
  166,68,26,153,199,37,123,58,100,134,216,91,5,231,
  185,140,210,48,110,237,179,81,15,78,16,242,172,47,
  113,147,205,17,79,173,243,112,46,204,146,211,141,
  111,49,178,236,14,80,175,241,19,77,206,144,114,44,
  109,51,209,143,12,82,176,238,50,108,142,208,83,13,
  239,177,240,174,76,18,145,207,45,115,202,148,118,
  40,171,245,23,73,8,86,180,234,105,55,213,139,87,9,
  235,181,54,104,138,212,149,203,41,119,244,170,72,
  22,233,183,85,11,136,214,52,106,43,117,151,201,74,
  20,246,168,116,42,200,150,21,75,169,247,182,232,10,
  84,215,137,107,53]

def crc8(data):
    c = 0
    for b in data:
        c = CRC8[c ^ b]
    return c

def pulse(angle):
    return int(500 + (max(0, min(180, angle)) / 180.0) * 2000)

def burst(ser, pan, tilt, dur_ms, m1, m2, m3, m4):
    pp, tp = pulse(pan), pulse(tilt)
    d = int(dur_ms)
    sd = bytearray([0x01, d & 0xFF, (d >> 8) & 0xFF, 2,
      2, pp & 0xFF, (pp >> 8) & 0xFF,
      1, tp & 0xFF, (tp >> 8) & 0xFF])
    fs = bytearray(b"\xaa\x55") + bytes([0x04, len(sd)]) + sd
    fs.append(crc8(fs[2:]))
    md = bytearray([0x05, 4])
    for mid, val in ((0, m1), (1, m2), (2, m3), (3, m4)):
        md += struct.pack("<Bf", mid, float(val))
    fm = bytearray(b"\xaa\x55") + bytes([0x03, len(md)]) + md
    fm.append(crc8(fm[2:]))
    ser.write(fs + fm)

ser = serial.Serial("/dev/ttyAMA0", 1_000_000)
burst(ser, pan=86, tilt=129, dur_ms=500, m1=200, m2=200, m3=0, m4=0)
ser.close()
\end{lstlisting}

\subsection{I\textsuperscript{2}C: leer ultras\'{o}nico + line follower}

\begin{lstlisting}[style=pystyle]
from smbus2 import SMBus, i2c_msg

bus = SMBus(1)

# --- Ultrasonico (0x77): leer distancia mm ---
w = i2c_msg.write(0x77, [0])
bus.i2c_rdwr(w)
r = i2c_msg.read(0x77, 2)
bus.i2c_rdwr(r)
dist_mm = int.from_bytes(bytes(list(r)), "little")
print(f"Distancia: {dist_mm} mm")

# --- Ultrasonico: LED verde solido ---
bus.write_i2c_block_data(0x77, 0x02, [0x00, 0x00, 0x10, 0x00,
                                       0x00, 0x10, 0x00])

# --- Line Follower (0x78): 4 sensores IR ---
v = bus.read_byte_data(0x78, 0x01)
sensors = [bool(v & b) for b in (0x01, 0x02, 0x04, 0x08)]
print(f"Sensores IR: {sensors}")
bus.close()
\end{lstlisting}

\subsection{C\'{a}mara: detecci\'{o}n HSV naranja}

\begin{lstlisting}[style=pystyle]
import cv2, numpy as np

cap = cv2.VideoCapture(0)
cap.set(cv2.CAP_PROP_FRAME_WIDTH, 320)
cap.set(cv2.CAP_PROP_FRAME_HEIGHT, 240)
cap.set(cv2.CAP_PROP_AUTO_EXPOSURE, 1)
cap.set(cv2.CAP_PROP_EXPOSURE, 200)

for _ in range(100):
    ok, frame = cap.read()
    if not ok:
        continue
    hsv = cv2.cvtColor(frame, cv2.COLOR_BGR2HSV)
    m1 = cv2.inRange(hsv, (0, 230, 20), (10, 255, 255))
    m2 = cv2.inRange(hsv, (168, 230, 20), (179, 255, 255))
    mask = cv2.bitwise_or(m1, m2)
    k = np.ones((3, 3), np.uint8)
    mask = cv2.dilate(cv2.morphologyEx(mask, cv2.MORPH_OPEN, k), k)
    contours, _ = cv2.findContours(mask, cv2.RETR_EXTERNAL,
                                    cv2.CHAIN_APPROX_SIMPLE)
    for c in contours:
        (x, y), r = cv2.minEnclosingCircle(c)
        if cv2.contourArea(c) > 700 and r > 15:
            print(f"Pelota ({int(x)}, {int(y)}) radio={r:.0f}")
            break
cap.release()
\end{lstlisting}

\newpage

% ============================================================
% EJEMPLOS RUST
% ============================================================

\section{Ejemplos Rust}

\textbf{Dependencias} (\texttt{Cargo.toml}):
\begin{lstlisting}[style=tomlstyle]
[dependencies]
serialport = "4"
i2cdev = "0.6"
opencv = { version = "0.93", features = ["clang-runtime"] }
\end{lstlisting}

\subsection{UART: enviar comando servo + motor}

\begin{lstlisting}[style=ruststyle]
use serialport::SerialPort;
use std::io::Write;

const CRC8: [u8; 256] = [0,94,188,226,97,63,221,131,
  194,156,126,32,163,253,31,65,157,195,33,127,252,
  162,64,30,95,1,227,189,62,96,130,220,35,125,159,
  193,66,28,254,160,225,191,93,3,128,222,60,98,190,
  224,2,92,223,129,99,61,124,34,192,158,29,67,161,
  255,70,24,250,164,39,121,155,197,132,218,56,102,
  229,187,89,7,219,133,103,57,186,228,6,88,25,71,
  165,251,120,38,196,154,101,59,217,135,4,90,184,
  230,167,249,27,69,198,152,122,36,248,166,68,26,
  153,199,37,123,58,100,134,216,91,5,231,185,140,
  210,48,110,237,179,81,15,78,16,242,172,47,113,
  147,205,17,79,173,243,112,46,204,146,211,141,
  111,49,178,236,14,80,175,241,19,77,206,144,114,
  44,109,51,209,143,12,82,176,238,50,108,142,208,
  83,13,239,177,240,174,76,18,145,207,45,115,202,
  148,118,40,171,245,23,73,8,86,180,234,105,55,
  213,139,87,9,235,181,54,104,138,212,149,203,41,
  119,244,170,72,22,233,183,85,11,136,214,52,106,
  43,117,151,201,74,20,246,168,116,42,200,150,21,
  75,169,247,182,232,10,84,215,137,107,53];

fn crc8(data: &[u8]) -> u8 {
    let mut c: u8 = 0;
    for &b in data { c = CRC8[(c ^ b) as usize]; }
    c
}

fn pulse(angle: u16) -> u16 {
    (500.0 + (angle.min(180) as f32 / 180.0) * 2000.0) as u16
}

fn main() -> Result<(), Box<dyn std::error::Error>> {
    let mut port = serialport::new("/dev/ttyAMA0", 1_000_000)
        .timeout(std::time::Duration::from_millis(100))
        .open()?;
    let (pan, tilt, dur) = (86u16, 129u16, 500u16);
    let (pp, tp) = (pulse(pan), pulse(tilt));
    let sd: Vec<u8> = vec![
        0x01, dur as u8, (dur >> 8) as u8, 2,
        2, pp as u8, (pp >> 8) as u8,
        1, tp as u8, (tp >> 8) as u8];
    let mut fs: Vec<u8> = vec![0xAA, 0x55, 0x04, sd.len() as u8];
    fs.extend(&sd);
    fs.push(crc8(&fs[2..]));
    let mut md: Vec<u8> = vec![0x05, 4];
    for (id, val) in [(0u8, 200f32), (1, 200.0), (2, 0.0), (3, 0.0)] {
        md.push(id);
        md.extend(&val.to_le_bytes());
    }
    let mut fm: Vec<u8> = vec![0xAA, 0x55, 0x03, md.len() as u8];
    fm.extend(&md);
    fm.push(crc8(&fm[2..]));
    port.write_all(&[&fs[..], &fm[..]].concat())?;
    println!("Trama enviada: {} bytes", fs.len() + fm.len());
    Ok(())
}
\end{lstlisting}

\subsection{I\textsuperscript{2}C: leer ultras\'{o}nico + line follower}

\begin{lstlisting}[style=ruststyle]
use i2cdev::core::*;
use i2cdev::linux::LinuxI2CDevice;

fn read_ultra(dev: &mut LinuxI2CDevice) -> Result<u16, ()> {
    dev.write(&[0x00]).map_err(|_| ())?;
    let mut buf = [0u8; 2];
    dev.read(&mut buf).map_err(|_| ())?;
    Ok(u16::from_le_bytes(buf).min(5000))
}

fn set_led(dev: &mut LinuxI2CDevice, r: u8, g: u8, b: u8)
    -> Result<(), ()> {
    dev.write(&[0x02, 0x00, r, g, b, r, g, b]).map_err(|_| ())
}

fn read_line(dev: &mut LinuxI2CDevice) -> Result<[bool; 4], ()> {
    dev.write(&[0x01]).map_err(|_| ())?;
    let mut buf = [0u8; 1];
    dev.read(&mut buf).map_err(|_| ())?;
    let v = buf[0];
    Ok([v & 1 != 0, v & 2 != 0, v & 4 != 0, v & 8 != 0])
}

fn main() -> Result<(), Box<dyn std::error::Error>> {
    let mut ultra = LinuxI2CDevice::new("/dev/i2c-1", 0x77)?;
    let mut line = LinuxI2CDevice::new("/dev/i2c-1", 0x78)?;
    let dist = read_ultra(&mut ultra)?;
    println!("Distancia: {} mm", dist);
    set_led(&mut ultra, 0x00, 0x10, 0x00)?;
    let sensors = read_line(&mut line)?;
    println!("Sensores IR: {:?}", sensors);
    Ok(())
}
\end{lstlisting}

\subsection{C\'{a}mara: detecci\'{o}n HSV naranja}

\begin{lstlisting}[style=ruststyle]
use opencv::prelude::*;
use opencv::{imgproc, videoio, core};

fn main() -> Result<(), Box<dyn std::error::Error>> {
    let mut cap = videoio::VideoCapture::new(0, videoio::CAP_ANY)?;
    cap.set(videoio::CAP_PROP_FRAME_WIDTH, 320.0)?;
    cap.set(videoio::CAP_PROP_FRAME_HEIGHT, 240.0)?;
    let mut frame = Mat::default();
    for _ in 0..100 {
        cap.read(&mut frame)?;
        if frame.empty() { continue; }
        let mut hsv = Mat::default();
        imgproc::cvt_color(&frame, &mut hsv,
            imgproc::COLOR_BGR2HSV, 0)?;
        let mut m1 = Mat::default();
        let mut m2 = Mat::default();
        core::in_range(&hsv,
            &core::Scalar::new(0.0, 230.0, 20.0, 0.0),
            &core::Scalar::new(10.0, 255.0, 255.0, 0.0),
            &mut m1)?;
        core::in_range(&hsv,
            &core::Scalar::new(168.0, 230.0, 20.0, 0.0),
            &core::Scalar::new(179.0, 255.0, 255.0, 0.0),
            &mut m2)?;
        let mut mask = Mat::default();
        core::bitwise_or(&m1, &m2, &mut mask,
            &core::no_array())?;
        let mut contours =
            opencv::types::VectorOfVectorOfPoint::new();
        imgproc::find_contours(&mut mask, &mut contours,
            imgproc::RETR_EXTERNAL,
            imgproc::CHAIN_APPROX_SIMPLE,
            core::Point::new(0, 0))?;
        for c in contours.iter() {
            let area = imgproc::contour_area(&c, false)?;
            let ((x, y), r) = imgproc::min_enclosing_circle(&c)?;
            if area > 700.0 && r > 15.0 {
                println!("Pelota ({:.0}, {:.0}) radio={:.0}", x, y, r);
                break;
            }
        }
    }
    Ok(())
}
\end{lstlisting}

\end{document}
